\documentclass[11pt,a4paper]{article}
\usepackage[a4paper, margin=2.2cm]{geometry}
\usepackage{graphicx}
\usepackage{booktabs}
\usepackage{xcolor}
\usepackage{listings}
\usepackage{hyperref}
\hypersetup{colorlinks=true, linkcolor=blue!60!black, urlcolor=blue!60!black}
\usepackage{amsmath, amssymb}
\usepackage{caption}
\usepackage{siunitx}
\sisetup{detect-all}
\usepackage{enumitem}
\usepackage{titling}
\usepackage{fancyhdr}
\usepackage[round,authoryear]{natbib}

\pagestyle{fancy}
\fancyhf{}
\rhead{\small Reproducibility Documentation}
\lhead{\small 4BS2 8-oxoG MD --- C. Ge, 2026}
\rfoot{\thepage}

% Code listing style
\definecolor{codegray}{rgb}{0.95,0.95,0.95}
\definecolor{codeblue}{rgb}{0.10,0.30,0.65}
\definecolor{codegreen}{rgb}{0.0,0.5,0.0}
\lstset{
  basicstyle=\ttfamily\footnotesize,
  backgroundcolor=\color{codegray},
  frame=single,
  rulecolor=\color{black!20},
  numbers=left,
  numberstyle=\tiny\color{black!50},
  keywordstyle=\color{codeblue}\bfseries,
  commentstyle=\color{codegreen}\itshape,
  stringstyle=\color{red!60!black},
  breaklines=true,
  postbreak=\mbox{\textcolor{red}{$\hookrightarrow$}\space},
  showstringspaces=false,
  tabsize=2,
}
\lstdefinelanguage{xml}{
  morekeywords={Residue,Atom,Bond,ExternalBond,HarmonicBondForce,Angle,
    HarmonicAngleForce,Proper,Improper,PeriodicTorsionForce,ForceField,Residues},
  morecomment=[s][\color{codegreen}]{<!--}{-->},
  sensitive=true,
}

\title{Reproducibility Documentation:\\
\large All-Atom MD Setup of 4BS2 with 8-Oxoguanosine RNA Modification}
\author{Caroline Ge\\\small carolinge73@gmail.com}
\date{April 30, 2026 (revision 2)}

\begin{document}
\maketitle

\begin{abstract}
\noindent
This document is a complete record of every step taken to set up,
parameterise, simulate, and analyse the wild-type 4BS2 protein--RNA
complex and its 8-oxoguanosine variant. It is written for an
investigator with general MD experience but no prior knowledge of this
project, with the explicit goal of bit-for-bit reproducibility.
Every input file, command line, parameter value, and software version
relevant to the published numbers is listed. All custom code is
tracked in the repository at \texttt{d:/all\_atom} (Windows host) and
deployed to \texttt{/data/biophys/carolinge/clawork/37\_OXR/} on the
PKS \texttt{newton} cluster (referred to as ``bio'' below).
\end{abstract}

\tableofcontents
\vspace{1em}

% =============================================================================
\section{Software environment}
\label{sec:env}

\subsection{Operating systems and hardware}

\begin{tabular}{ll}
\toprule
Workstation              & Windows 10/11, Python 3.11, used for setup, analysis, plotting \\
Cluster                  & \texttt{newton.pks.mpg.de}, Slurm scheduler (CentOS-style RHEL) \\
GPUs                     & NVIDIA A100 (mixed precision; CUDA Toolkit 12.5 module) \\
Cluster home             & \texttt{/home/carolinge}; \textbf{80 GB quota} \\
Cluster scratch          & \texttt{/data/biophys/carolinge/} (1.4 PB pool, 444 TB free) \\
\bottomrule
\end{tabular}

\paragraph{Important note on disk layout.}
The cluster home filesystem (\texttt{/home}) has an 80 GB user quota.
A single 100\,ns trajectory at the parameters used here is $\sim$3 GB,
so all production output must be written to \texttt{/data}. We use a
symlink to keep paths consistent:
\begin{lstlisting}[language=bash]
ln -s /data/biophys/carolinge/clawork/37_OXR ~/work/37_OXR
\end{lstlisting}

\subsection{Conda environment}

The environment name on the cluster is \texttt{allatom}. Pinned versions
of the load-bearing packages (\texttt{conda list}):

\begin{tabular}{ll}
\toprule
\texttt{python}               & 3.11 \\
\texttt{openmm}               & 8.2.0 \\
\texttt{cuda-version}         & 12.5 \\
\texttt{cuda-nvrtc}           & 12.5.82 \\
\texttt{pdbfixer}             & 1.10 \\
\texttt{mdanalysis}           & 2.10.0 \\
\texttt{numpy}                & 2.x \\
\texttt{biopython}            & 1.85 \\
\texttt{matplotlib}           & 3.9.x \\
\bottomrule
\end{tabular}

A frozen environment specification is committed to
\texttt{environment\_server.yml} in the repository root.

% =============================================================================
\section{Project layout}
\label{sec:layout}

\subsection{Workstation (\texttt{d:/all\_atom/})}

\begin{lstlisting}[language=bash]
d:/all_atom/
  app/
    server.py                    # Flask UI for the MD setup pipeline
    ptm_builder.py               # NeRF-based PTM applicator (defines 8OG patch)
    ff_custom/
      8OG_RNA_amber.xml          # Custom force field for 8-oxoguanosine
    templates/index.html         # Web UI
  scripts/
    gen_4bs2_md.py               # Generate WT and 8OG starting structures
    gen_replicas.py              # Generate 3 WT + 3 8OG replica directories
    analyze_binding.py           # Single-trajectory binding analysis
    analyze_rmsf.py              # Single-trajectory per-residue RMSF
    analyze_replicas.py          # Multi-replica analysis with partial-cov averaging
    plot_rmsf.py                 # Plot RMSF comparison
    compare_wt_8og.py            # Plot WT vs 8OG comparison panels
  output/
    4BS2_WT_MD/                  # 10 ns preliminary WT run
    4BS2_8OG_G3_MD/              # 10 ns preliminary 8OG run
    4BS2_WT_r{1,2,3}/            # Local staging of 100 ns replica directories
    4BS2_8OG_r{1,2,3}/
    analysis/                    # All analysis CSVs and PNGs
  docs/
    4BS2_8OG_report.tex          # Main results report
    4BS2_8OG_methods.tex         # This file
\end{lstlisting}

\subsection{Cluster (\texttt{/data/biophys/carolinge/clawork/37\_OXR/})}

\begin{lstlisting}[language=bash]
/data/biophys/carolinge/clawork/37_OXR/
  04_4BS2_WT_MD/                 # 10 ns preliminary WT run
  05_4BS2_8OG_G3_MD/             # 10 ns preliminary 8OG run
  replicas/
    4BS2_WT_r1/   ...   _r3/     # 100 ns WT replicas (3 GB DCD each)
    4BS2_8OG_r1/  ...   _r3/     # 100 ns 8OG replicas
  anal_gyrate/                   # All analysis output: CSVs, NPZs, PNGs
  analyze_binding.py             # uploaded copies of the local scripts
  analyze_rmsf.py
  analyze_replicas.py
\end{lstlisting}

% =============================================================================
\section{Source structure: 4BS2}
\label{sec:source}

The structure used as the starting point is PDB entry \textbf{4BS2}
(\citealt{Lukavsky2013}; an RRM--RNA complex). The CIF file
(\texttt{d:/all\_atom/4BS2.cif}) was downloaded from the RCSB PDB and
contains:

\begin{itemize}[leftmargin=*]
\item Chain A: 174 amino acids (the RRM domain, GLY1--GLN174);
\item Chain B: 12 nucleotides of single-stranded RNA with sequence
\texttt{GUGUAGAAUGAAU}.
\end{itemize}

The third nucleotide on chain B (G3) is the target of the 8-oxo
modification.

% =============================================================================
\section{Force field for 8-oxoguanosine}
\label{sec:ff}

The wild-type RNA force field is AMBER OL3 (\citealt{Perez2007},
\texttt{amber14/RNA.OL3.xml} as shipped with OpenMM~8.2). 8OG requires
a custom XML (\texttt{app/ff\_custom/8OG\_RNA\_amber.xml}) that
\emph{augments}, not replaces, OL3. The two XMLs are loaded together
when the system is built (Section~\ref{sec:run}).

\subsection{Atomic changes vs. canonical guanosine}

\begin{tabular}{lll}
\toprule
Atom & Change & Notes \\
\midrule
H8 & deleted & aromatic CH on C8 in G is removed \\
C8 & atom type RNA-CP $\rightarrow$ RNA-C & sp$^{2}$ aromatic $\rightarrow$ sp$^{2}$ carbonyl C \\
N7 & atom type RNA-NB $\rightarrow$ RNA-NA & aromatic N $\rightarrow$ sp$^{2}$ N--H \\
O8 & added (RNA-O) & C8=O8 carbonyl oxygen \\
H7 & added (RNA-H) & N7--H proton (donor) \\
\bottomrule
\end{tabular}

No new atom types are introduced: RNA-C, RNA-O, RNA-NA, RNA-H, RNA-CB,
RNA-N* are all defined in \texttt{amber14/RNA.OL3.xml}. Only previously
unused \emph{combinations} of these types appear in 8OG, which is why
new bond/angle/torsion entries are needed (Section~\ref{sec:ff:supp}).

\subsection{Partial charges}
\label{sec:ff:charges}

Backbone and sugar charges are copied verbatim from the standard G
residue in OL3.xml. Base charges are unchanged from G except in the
N7--C8--O8--H7 region. Modified charges (in elementary charge units):

\begin{tabular}{lcc}
\toprule
Atom & G (OL3) & 8OG (this work) \\
\midrule
C8       & $+0.1374$  (RNA-CP) & $+0.4000$ (RNA-C, carbonyl) \\
H8       & $+0.1640$  (deleted) & --- \\
N7       & $-0.5709$  (RNA-NB) & $-0.4600$ (RNA-NA) \\
H7       & ---                  & $+0.3500$ (RNA-H) \\
O8       & ---                  & $-0.5600$ (RNA-O) \\
\bottomrule
\end{tabular}

\paragraph{Reference and limitations.} Values are derived by analogy
from \citet{Aduri2007}, who report RESP charges for several modified
RNA nucleosides. The exact 8OG charges in their work are for the
DNA-context 8OG; for the RNA context we kept the OL3 backbone (which
already differs from the Aduri scheme) and only redistributed charge
locally around C8/N7. \textbf{Net residue charge is preserved at
$-1.0$} (the same as canonical phosphorylated G):
\begin{equation}
\Delta q_{\text{base}} =
   (-0.1640) + (+0.2626) + (+0.1109) + (+0.3500) + (-0.5600) = -0.0005 \approx 0.
\end{equation}

For publication-quality work this charge model should be replaced with
a fresh QM RESP fit using the Antechamber pipeline. The validation
script (\texttt{scripts/validate\_8OG\_ff.py}) confirms that the
present charges and topology load cleanly into OpenMM and produce a
finite-energy system, but does not validate accuracy.

\subsection{Supplemental bonded parameters}
\label{sec:ff:supp}

OL3 lacks parameters for the type-pair \texttt{RNA-CB}--\texttt{RNA-NA}
because canonical guanine has \texttt{RNA-NB} at N7. Five bonded terms
are added by analogy to existing OL3 entries:

\begin{lstlisting}[language=xml,caption={Supplemental bonded parameters in
\texttt{8OG\_RNA\_amber.xml}},captionpos=b]
<HarmonicBondForce>
  <!-- C5(CB)-N7(NA) bond, derived from CB-NB in OL3.xml: k=346435, r=0.1391 nm -->
  <Bond k="346435.2" length="0.1390"
        type1="RNA-CB" type2="RNA-NA"/>
</HarmonicBondForce>

<HarmonicAngleForce>
  <!-- C8-N7-C5 angle in 5-membered imidazole ring (~103.8 deg = 1.8117 rad)
       analogous to RNA-CB-RNA-NB-RNA-CP in OL3 -->
  <Angle angle="1.8117" k="585.76"
         type1="RNA-CB" type2="RNA-NA" type3="RNA-C"/>

  <!-- C5-N7-H7 exocyclic angle (~139.4 deg = 2.433 rad).
       Derived from planarity: 360 - 103.8 (C8-N7-C5) - 116.8 (C8-N7-H7) -->
  <Angle angle="2.433" k="418.40"
         type1="RNA-CB" type2="RNA-NA" type3="RNA-H"/>
</HarmonicAngleForce>

<PeriodicTorsionForce ordering="amber">
  <!-- X-CB-NA-X torsion, analogous to X-CB-NB-X in OL3 -->
  <Proper k1="10.6692" periodicity1="2" phase1="3.141592653589793"
          type1="" type2="RNA-CB" type3="RNA-NA" type4=""/>

  <!-- N* improper to enforce planarity at N9 (analogous to G in OL3,
       updated for the new C8 atom type) -->
  <Improper k1="4.184" periodicity1="2" phase1="3.141592653589793"
            type1="RNA-N*" type2="RNA-CB" type3="RNA-C" type4="RNA-CT"/>
</PeriodicTorsionForce>
\end{lstlisting}

All other 8OG bonded interactions use parameters that already exist in
OL3 for canonical RNA atom-type pairs.

\subsection{Residue templates}
Three residue templates are defined to handle the three positional
contexts of 8OG in an RNA chain:
\begin{itemize}[leftmargin=*]
\item \texttt{8OG} -- internal residue, \texttt{ExternalBond} on
both \texttt{P} (5$'$ side) and \texttt{O3'} (3$'$ side).
\item \texttt{8OG5} -- 5$'$-terminal: no phosphate; \texttt{HO5'} caps
\texttt{O5'}; only \texttt{O3'} forms an external bond.
\item \texttt{8OG3} -- 3$'$-terminal: no external bond on \texttt{O3'};
\texttt{HO3'} cap added.
\end{itemize}

\subsection{Validation}
The script \texttt{scripts/validate\_8OG\_ff.py} loads OL3 + the custom
XML, builds a 5$'$-terminal 8OG mononucleotide topology, and calls
\texttt{ff.createSystem()}. Failure of this step (for example, after
introducing a typo in the XML) indicates that the parameter set is
incomplete; success indicates that \emph{at least} the bonded topology
is consistent.

% =============================================================================
\section{Initial structure preparation}
\label{sec:prep}

The full preparation pipeline is encapsulated in the script
\texttt{scripts/gen\_4bs2\_md.py}; the key steps are summarised here so
the reader can recreate them by hand if needed.

\subsection{Wild-type structure}
\begin{enumerate}[leftmargin=*]
\item Load \texttt{4BS2.cif} via PDBFixer.
\item Call \texttt{findMissingResidues()} (no terminal trimming).
\item Call \texttt{findMissingAtoms()} and \texttt{addMissingAtoms()}.
\item Call \texttt{addMissingHydrogens(pH=7.0)}.
\item Write the result to \texttt{4BS2\_prepared.pdb}
(3164 heavy + H atoms, single MODEL).
\end{enumerate}

\subsection{8OG variant: residue surgery via NeRF}
\label{sec:prep:nerf}

The 8OG variant is produced from the prepared WT PDB by an idealised
geometric procedure implemented in
\texttt{app/ptm\_builder.py}. The procedure for chain B residue 3 is:

\begin{enumerate}[leftmargin=*]
\item Read the WT PDB with BioPython.
\item Locate the target residue (chain B, resseq 3, resname G).
\item Delete atom H8 from the residue.
\item Generate the position of \texttt{O8} using the Natural Extension
Reference Frame (NeRF) algorithm with three reference atoms (N9, N7, C8)
and ideal internal coordinates:
\begin{itemize}[leftmargin=*]
\item Bond \texttt{C8--O8}: 1.230\,\si{\angstrom}
\item Angle \texttt{N7--C8--O8}: 128.8\textdegree
\item Dihedral \texttt{N9--N7--C8--O8}: 180\textdegree (O8 trans to N9)
\end{itemize}
\item Generate the position of \texttt{H7} from references (C8, C5, N7):
\begin{itemize}[leftmargin=*]
\item Bond \texttt{N7--H7}: 1.010\,\si{\angstrom}
\item Angle \texttt{C5--N7--H7}: 125.2\textdegree
\item Dihedral \texttt{C8--C5--N7--H7}: 180\textdegree (H7 in imidazole plane)
\end{itemize}
\item Rename the residue to \texttt{8OG} (or \texttt{8OG3}/\texttt{8OG5}
if it is the 3$'$ or 5$'$ terminus).
\item Write the modified structure to PDB.
\item Append \texttt{CONECT} records for all intra-residue bonds of 8OG
\emph{and} for the inter-residue O3$'$--P backbone bonds spanning the
modified residue. This is essential because OpenMM's
\texttt{createStandardBonds()} routine builds bonds only for residues
named in its built-in standard table; for any non-standard name it
silently skips both the residue's internal topology \emph{and} its
inter-residue backbone bonds. Without the explicit CONECT records,
\texttt{ForceField.createSystem()} fails with the misleading message
``\emph{atoms match 8OG, but the bonds are different}''.
\end{enumerate}

\paragraph{NeRF algorithm.}
The NeRF subroutine takes three known atom positions $\mathbf{a},
\mathbf{b}, \mathbf{c}$ and three internal coordinates (bond length $b$,
bond angle $\theta$, dihedral $\phi$), and returns the position of a
new atom $\mathbf{d}$ such that
$|\mathbf{c}-\mathbf{d}| = b$,
$\angle(\mathbf{b}, \mathbf{c}, \mathbf{d}) = \theta$, and
$\text{dihedral}(\mathbf{a}, \mathbf{b}, \mathbf{c}, \mathbf{d}) = \phi$.
The implementation in \texttt{ptm\_builder.py} is a 7-line numpy
function that constructs the local frame at $\mathbf{c}$.

\paragraph{Reproduction.}
\begin{lstlisting}[language=bash]
python -c "
import sys; sys.path.insert(0, 'd:/all_atom/app')
from ptm_builder import apply_ptm
pdb = apply_ptm('output/4BS2_WT_MD/4BS2_prepared.pdb',
                 chain_id='B', resnum=3, ptm_code='8OG')
open('output/4BS2_8OG_G3_MD/4BS2_8OG_G3.pdb','w').write(pdb)
"
\end{lstlisting}

\subsection{File outputs}
After preparation, two files are used as the starting structure of all
subsequent simulations:

\begin{tabular}{ll}
\toprule
\texttt{output/4BS2\_WT\_MD/4BS2\_prepared.pdb} & 3164 atoms; canonical RRM--RNA \\
\texttt{output/4BS2\_8OG\_G3\_MD/4BS2\_8OG\_G3.pdb} & 3165 atoms; G3 $\rightarrow$ 8OG (--H8 +O8 +H7) \\
\bottomrule
\end{tabular}

% =============================================================================
\section{Simulation protocol}
\label{sec:run}

The simulation protocol is identical for the preliminary 10\,ns runs
(Section~\ref{sec:run:short}) and the production 100\,ns multi-replica
runs (Section~\ref{sec:run:long}); only the production length and the
random seed differ.

\subsection{Force-field stack}
\begin{lstlisting}[language=python]
ff = ForceField('amber14/protein.ff14SB.xml',     # protein
                'amber14/RNA.OL3.xml',             # RNA
                'amber14/tip3p.xml',               # water
                'app/ff_custom/8OG_RNA_amber.xml') # 8OG (only for variant)
\end{lstlisting}

\subsection{Solvation}
\texttt{Modeller.addSolvent} with:
\begin{itemize}[leftmargin=*]
\item Padding: 1.0\,nm (cubic box; box size scales with structure)
\item Ionic strength: 0.15\,M
\item Cation: Na$^{+}$; anion: Cl$^{-}$
\item Total system: $\sim$51\,407 atoms (WT) / $\sim$51\,408 (8OG, +1 atom)
\end{itemize}

\subsection{Energy minimisation and equilibration}
\begin{enumerate}[leftmargin=*]
\item Energy minimisation: \texttt{minimizeEnergy(maxIterations=1000)}.
\item NVT equilibration: 100\,ps.
LangevinMiddleIntegrator (300\,K, $\gamma = 1$\,ps$^{-1}$, dt =
4\,fs, hydrogen mass repartitioning to 1.5\,amu).
PME with 1.2\,nm cutoff. HBonds constraints. State logged every 500
steps to \texttt{nvt.log}.
\item NPT equilibration: 100\,ps. As above, plus
MonteCarloBarostat (1\,atm, 300\,K). State logged every 500 steps to
\texttt{npt.log}.
\end{enumerate}

\subsection{Production}
\label{sec:run:long}
\begin{itemize}[leftmargin=*]
\item Length: 100\,ns (25\,000\,000 steps at 4\,fs).
\item Trajectory: \texttt{DCDReporter('traj.dcd', 5000)}, i.e.\ a frame
every 20\,ps, $\sim$5000 frames per replica, $\sim$3.0 GB per file.
\item State log: every 5000 steps to \texttt{md.log}.
\item Final checkpoint saved to \texttt{final.chk}.
\item Hardware: NVIDIA A100, mixed precision; speed $\sim$470\,ns/day,
wall-clock 5\,h per 100\,ns replica.
\end{itemize}

\subsection{Replica generation}
The script \texttt{scripts/gen\_replicas.py} generates 6 directories
(3 WT + 3 8OG), each with its own \texttt{run\_openmm.py} and
\texttt{sub.sh}. Replicas differ in the integer \emph{seed} ($1, 2,
\dots, 6$) which is used at three places to ensure independent samples:

\begin{lstlisting}[language=python]
random.seed(seed)                                  # Python RNG, used by Modeller
integrator.setRandomNumberSeed(seed)               # Langevin thermostat noise
sim.context.setVelocitiesToTemperature(300*K, seed) # initial velocities
barostat.setRandomNumberSeed(seed)                 # MC barostat
\end{lstlisting}

\subsection{Slurm submission script}
\begin{lstlisting}[language=bash,caption={\texttt{sub.sh} (one per
replica)}]
#!/bin/bash
#SBATCH --time=24:00:00
#SBATCH --partition=graphic
#SBATCH --constraint="gpu"
#SBATCH --gres=gpu:a100:1
#SBATCH --cpus-per-task=10
#SBATCH --mem=225000
#SBATCH --job-name="4BS2_<system>_r<n>"
#SBATCH --output=./%j.out
#SBATCH --error=./%j.err

set -e
module purge
module load cuda/12.5
source "$(conda info --base)/etc/profile.d/conda.sh"
conda activate allatom
python ./run_openmm.py
\end{lstlisting}

To submit one replica:
\begin{lstlisting}[language=bash]
cd /data/biophys/carolinge/clawork/37_OXR/replicas/4BS2_WT_r1
sbatch sub.sh
\end{lstlisting}

\subsection{Output files per replica}
After a successful run, the directory contains:
\begin{tabular}{lp{8.6cm}}
\toprule
\texttt{4BS2\_<sys>.pdb} & starting structure (input) \\
\texttt{8OG\_RNA\_amber.xml} & custom FF (only for 8OG variant) \\
\texttt{run\_openmm.py} & simulation driver (input) \\
\texttt{sub.sh} & submission script (input) \\
\texttt{solvated.pdb} & solvated topology written before MD; matches DCD atom order; required for analysis \\
\texttt{nvt.log}, \texttt{npt.log} & equilibration state logs \\
\texttt{md.log} & production state log; contains progress \% used by analysis script \\
\texttt{traj.dcd} & production trajectory (5000 frames, $\sim$3 GB) \\
\texttt{final.chk} & checkpoint at end of production \\
\bottomrule
\end{tabular}

\subsection{Resuming a failed replica}
A replica that died mid-production (for example because of a
filesystem-quota error) cannot be resumed mid-run because no
intermediate checkpoint is written. Restart from scratch:

\begin{lstlisting}[language=bash]
cd /data/biophys/carolinge/clawork/37_OXR/replicas/4BS2_8OG_r3
rm -f traj.dcd md.log nvt.log npt.log solvated.pdb final.chk *.out *.err
sbatch sub.sh
\end{lstlisting}

% =============================================================================
\section{Analysis}
\label{sec:anal}

All analysis scripts are kept under \texttt{scripts/} on the
workstation and are deployed to the cluster's project root for
execution against trajectories on \texttt{/data}.

\subsection{Single-trajectory binding analysis}
\texttt{scripts/analyze\_binding.py} computes, on a per-frame basis:
RNA RMSD vs.\ frame 0; total heavy-atom contacts ($<4\,\si{\angstrom}$)
between protein and RNA; protein--RNA centre-of-mass distance;
G3-specific contacts; G3 base polar contacts ($<3.5\,\si{\angstrom}$,
proxy for hydrogen bonds). It also accumulates a per-residue
contact-frequency map. Inputs: a topology PDB (must match the DCD atom
count) and a trajectory file. Outputs: \texttt{<prefix>\_cv.csv,
\_contact\_map.csv, \_summary.txt} and PNG plots.

The topology PDB must be the \texttt{solvated.pdb} written by
\texttt{run\_openmm.py}, not the unsolvated starting file --- the DCD
contains all 51\,407+ atoms.

\subsection{Single-trajectory RMSF}
\texttt{scripts/analyze\_rmsf.py} aligns on protein C$\alpha$ and
computes per-residue heavy-atom RMSF for both chains.

\subsection{Multi-replica analysis with partial-coverage averaging}
\label{sec:anal:multi}
\texttt{scripts/analyze\_replicas.py} is the main workhorse for the
100\,ns multi-replica analysis. Key features:

\begin{itemize}[leftmargin=*]
\item Accepts an arbitrary number of \texttt{--label / --replicas}
pairs.
\item Loads each replica's \texttt{solvated.pdb} + \texttt{traj.dcd}
independently. Trajectories of unequal length are handled by NaN-padding
when stacked, so that at each time-point the cross-replica mean and SEM
are computed only over replicas that have actually reached that
time-point.
\item Applies a \emph{minimum-image PBC correction} relative to the
protein centre of mass before computing RNA RMSF and RMSD; this is
necessary because RNA atoms can wander to a periodic image, which would
otherwise inflate RMSF artifactually.
\item Output files: \texttt{multi\_per\_replica.json} (per-replica means),
\texttt{multi\_timeseries.npz/.png}, \texttt{multi\_rmsf.npz/.png}, and
\texttt{multi\_summary.txt}.
\end{itemize}

\paragraph{Avoiding the active-trajectory race condition.}
Reading a DCD file that is concurrently being written by an active MD
job will cause the analysis script to die silently mid-frame. The
analysis must therefore exclude any replica whose simulation has not
yet finished. A simple check is whether the corresponding Slurm job is
still in the queue:
\begin{lstlisting}[language=bash]
squeue -u carolinge | grep 4BS2
\end{lstlisting}
Only fully completed replicas should be passed to
\texttt{analyze\_replicas.py}.

\paragraph{Reproduction.}
\begin{lstlisting}[language=bash]
cd /data/biophys/carolinge/clawork/37_OXR
nohup conda run -n allatom python analyze_replicas.py \
  --label WT  --replicas replicas/4BS2_WT_r1 \
                          replicas/4BS2_WT_r2 \
                          replicas/4BS2_WT_r3 \
  --label 8OG --replicas replicas/4BS2_8OG_r1 \
                          replicas/4BS2_8OG_r2 \
                          replicas/4BS2_8OG_r3 \
  --out anal_gyrate/multi --stride 5 \
  > /tmp/multi.log 2>&1 &
\end{lstlisting}

The \texttt{--stride 5} option subsamples every 5th frame from each
trajectory (1000 effective frames per replica) and reduces wall time
to $\sim$5 minutes for 6 replicas.

\subsection{Plot generation}
\texttt{scripts/compare\_wt\_8og.py} (10\,ns single-replica figures) and
\texttt{scripts/plot\_rmsf.py} (single-replica RMSF) read the CSVs/JSON
emitted by the analysis scripts and write the corresponding PNGs into
\texttt{output/analysis/}.

% =============================================================================
\section{Step-by-step reproduction recipe}
\label{sec:repro}

This section is a single-command-per-step recipe to recreate every
file from scratch.

\subsection{One-time setup (workstation)}

\begin{lstlisting}[language=bash]
# 1. Clone repository (already in place at d:/all_atom)
# 2. Create conda environment
conda env create -f environment_server.yml -n allatom

# 3. Verify the custom force field loads
conda run -n allatom python scripts/validate_8OG_ff.py
\end{lstlisting}

\subsection{One-time setup (cluster)}

\begin{lstlisting}[language=bash]
ssh bio
# 4. Mirror the conda environment
conda env create -f environment_server.yml -n allatom
mkdir -p /data/biophys/carolinge/clawork/37_OXR
ln -s /data/biophys/carolinge/clawork/37_OXR ~/work/37_OXR
\end{lstlisting}

\subsection{Generate starting structures}

\begin{lstlisting}[language=bash]
# Workstation
cd d:/all_atom
conda run -n allatom python scripts/gen_4bs2_md.py
# Outputs:
#   output/4BS2_WT_MD/4BS2_prepared.pdb
#   output/4BS2_8OG_G3_MD/4BS2_8OG_G3.pdb (CONECT records included)
\end{lstlisting}

\subsection{Generate replica directories}

\begin{lstlisting}[language=bash]
conda run -n allatom python scripts/gen_replicas.py
# Outputs 6 directories under output/4BS2_{WT,8OG}_r{1,2,3}/
\end{lstlisting}

\subsection{Upload to cluster and submit}

\begin{lstlisting}[language=bash]
# Workstation -> cluster
scp -r output/4BS2_*_r? bio:/data/biophys/carolinge/clawork/37_OXR/replicas/
ssh bio
cd /data/biophys/carolinge/clawork/37_OXR/replicas
for d in 4BS2_*_r*; do
  dos2unix $d/sub.sh $d/run_openmm.py 2>/dev/null
  ( cd $d && sbatch sub.sh )
done
squeue -u $USER
\end{lstlisting}

\subsection{Wait, then analyse}

\begin{lstlisting}[language=bash]
# Wait until squeue -u $USER returns no 4BS2_* jobs.
cd /data/biophys/carolinge/clawork/37_OXR
nohup conda run -n allatom python analyze_replicas.py \
  --label WT  --replicas replicas/4BS2_WT_r{1,2,3} \
  --label 8OG --replicas replicas/4BS2_8OG_r{1,2,3} \
  --out anal_gyrate/multi --stride 5 \
  > /tmp/multi.log 2>&1 &

# Download outputs to workstation:
scp 'bio:/data/biophys/carolinge/clawork/37_OXR/anal_gyrate/multi_*' \
    d:/all_atom/output/analysis/

# Compile the report
cd d:/all_atom/docs
pdflatex -interaction=nonstopmode 4BS2_8OG_report.tex
pdflatex -interaction=nonstopmode 4BS2_8OG_report.tex   # citations
\end{lstlisting}

% =============================================================================
\section{Known issues and gotchas}
\label{sec:gotchas}

\begin{enumerate}[leftmargin=*]
\item \textbf{Disk quota on \texttt{/home}.}
80 GB user quota; a single 100 ns trajectory is $\sim$3 GB. All
production output must live on \texttt{/data}. Failure to do so
manifests as a Slurm job that exits without an error message after
$\sim$40\,\% of production. The symlink trick in Section~\ref{sec:env}
keeps the path \texttt{\textasciitilde/work/37\_OXR} usable while
storing data on \texttt{/data}.

\item \textbf{Race condition on active trajectories.}
\texttt{analyze\_replicas.py} will silently die mid-frame if it reads
a DCD that is being concurrently written by an active MD job. Always
verify with \texttt{squeue} that all replicas are complete before
running the multi-replica analysis. (The script uses NaN-padding to
handle replicas of \emph{differing} lengths, but only after they have
finished writing.)

\item \textbf{CONECT records are essential for non-standard residues.}
OpenMM's \texttt{createStandardBonds()} bond-detection routine skips
any residue whose name is not in its built-in standard table. For 8OG
this means \emph{none} of the residue's intra-residue bonds will be
created, and the inter-residue O3$'$--P backbone bonds spanning the
modified residue will also be silently dropped. The result is an
opaque error message at \texttt{ForceField.createSystem()} time:
\emph{``atoms match 8OG, but the bonds are different''}. The
\texttt{ptm\_builder.py} script writes explicit CONECT records for
both the intra-residue bonds (Section~\ref{sec:prep:nerf}) and the
inter-residue backbone bonds (function
\texttt{\_add\_backbone\_conect}). Do not modify the output PDB without
preserving these CONECT records.

\item \textbf{Windows-style line endings break Slurm.}
Submission scripts edited or copied via Windows utilities can pick up
CRLF line endings, which cause sbatch to reject the script with a
``DOS line breaks'' error. Run \texttt{dos2unix} on every \texttt{*.sh}
and \texttt{*.py} after upload.

\item \textbf{CUDA driver / PTX version mismatch.}
\texttt{\$ nvidia-smi} on the login node may report a different driver
version than what is installed on the GPU compute nodes. If
\texttt{Platform.getPlatformByName('CUDA')} succeeds but a subsequent
\texttt{Simulation()} call fails with ``CUDA\_ERROR\_UNSUPPORTED\_PTX\_VERSION'',
fall back to OpenCL or CPU automatically (the
\texttt{run\_openmm.py} template includes such a fallback loop).

\item \textbf{PBC unwrapping for RNA RMSF.}
A naive RMSF computation will report values of $\sim$25--40\,\si{\angstrom}
for any RNA atoms that hop to a periodic image. \texttt{analyze\_replicas.py}
applies a minimum-image correction relative to the protein centre of mass
before storing positions. Without this correction, RNA RMSF is
uninterpretable.
\end{enumerate}

% =============================================================================
\renewcommand*{\bibfont}{\small}
\begin{thebibliography}{9}

\bibitem[Aduri et~al.(2007)]{Aduri2007}
Aduri R., Psciuk B.T., Saro P., Taniga H., Schlegel H.B., SantaLucia J.
\newblock {AMBER force field parameters for the naturally occurring modified
nucleosides in RNA}.
\newblock \emph{J. Chem. Theory Comput.} \textbf{3}, 1464--1475 (2007).

\bibitem[Lukavsky et~al.(2013)]{Lukavsky2013}
Lukavsky P.J., Daujotyte D., Tollervey J.R., Ule J., Stuani C., Buratti E.,
Baralle F.E., Damberger F.F., Allain F.H.-T.
\newblock {Molecular basis of UG-rich RNA recognition by the human splicing
factor TDP-43}.
\newblock \emph{Nat. Struct. Mol. Biol.} \textbf{20}, 1443--1449 (2013).

\bibitem[P\'erez et~al.(2007)]{Perez2007}
P\'erez A., March\'an I., Svozil D., \v{S}poner J., Cheatham T.E.,
Laughton C.A., Orozco M.
\newblock {Refinement of the AMBER force field for nucleic acids:
improving the description of $\alpha$/$\gamma$ conformers}.
\newblock \emph{Biophys. J.} \textbf{92}, 3817--3829 (2007).

\end{thebibliography}

\end{document}
